\chapter{The Canonical Dihedral Catalogue}
\label{chap:canonical-taxonomy}

\section{Taxonomic Ordering Principles}

The $509$ dihedral composition mosaics across $\mathbb{Z}_1$ through $\mathbb{Z}_8$ are organised according to a four-tier hierarchical classification:

\begin{enumerate}
    \item \textbf{Tier 1 (Lattice Ring Progression)}: Mosaics are stratified primarily by the underlying circular lattice order $n \in \{1, 2, \dots, 8\}$.
    
    \item \textbf{Tier 2 (Part-Count Stratification)}: Within each ring $\mathbb{Z}_n$, mosaics are stratified by the number of constituent blocks $s = \ell(\lambda) \in \{1, \dots, n\}$, corresponding to the Stirling numbers of the second kind.
    
    \item \textbf{Tier 3 (Lexicographical Shape Ordering)}: Within each stratum $s$, individual partition shapes $\lambda = (\lambda_1, \dots, \lambda_s) \vdash n$ follow standard lexicographical order:
    \begin{equation}
    (n) \prec (n-1, 1) \prec (n-2, 2) \prec \dots \prec (2, 1^{n-2}) \prec (1^n)
    \end{equation}
    
    \item \textbf{Tier 4 (Intra-Shape Ordering)}: Within each individual shape $\lambda$, non-equivalent dihedral orbits are ordered deterministically by:
    \begin{itemize}
        \item \textbf{Symmetric Block Count ($s_{\text{sym}}$ Descending)}: Fully achiral configurations precede mixed configurations, which precede chiral configurations.
        \item \textbf{Chord Energy ($E$ Descending)}: Mosaics spanning longer circular distances precede localized configurations:
        \begin{equation}
        E(\mathcal{M}) = \sum_{B \in \mathcal{M}} \sum_{u < v \in B} \min(|u - v|, n - |u - v|)^2
        \end{equation}
        \item \textbf{Canonical Representative String}: Deterministic tie-breaking is governed by the lexicographical order of the hexadecimal slash-delimited partition notation.
    \end{itemize}
\end{enumerate}

Every orbit is assigned both a global integer identifier $1 \dots 509$ spanning the full octonary universe and a local ring identifier $1 \dots M_n$, directly matching its appearance in the visual plates of Chapter~\ref{chap:visual-atlas}.

\section{Stratified Partition Shape Inventory}

Table~\ref{tab:shape-census-z1-z8} details the complete partition shape inventory across all eight rings $\mathbb{Z}_1$ through $\mathbb{Z}_8$, recording the number of set partitions $p_\lambda$, the number of non-equivalent dihedral mosaics, the shape monomial $\mathbf{z}^\lambda$, and the bivariate orbit polynomial $Q_\lambda(\omega)$.

% Auto-generated by generate_multilattice_atlas.py
\begin{longtable}{r r l l}
\caption{\label{tab:shape-census-z1-z8}Stratified Partition Shape Inventory for $\mathbb{Z}_1$ through $\mathbb{Z}_8$}\\
\toprule
Partitions $\#p$ & Mosaics $\#o$ & Partition Shape Monomial $\mathbf{z}^\lambda$ & Orbit Polynomial $Q_\lambda(\omega)$ \\
\midrule
\endfirsthead
\caption[]{Stratified Partition Shape Inventory (continued)}\\
\toprule
Partitions $\#p$ & Mosaics $\#o$ & Partition Shape Monomial $\mathbf{z}^\lambda$ & Orbit Polynomial $Q_\lambda(\omega)$ \\
\midrule
\endhead
\bottomrule
\endfoot
\midrule
\multicolumn{4}{c}{\textbf{Ring $\mathbb{Z}_1$} ($|D_1|=1$, $B_1=1$, 1 Mosaics)} \\
\midrule
1 & 1 & $z_1$ & $1\omega^{1}$ \\
\midrule
\multicolumn{4}{c}{\textbf{Ring $\mathbb{Z}_2$} ($|D_2|=2$, $B_2=2$, 2 Mosaics)} \\
\midrule
1 & 1 & $z_2$ & $1\omega^{1}$ \\
1 & 1 & $z_1^2$ & $1\omega^{1}$ \\
\midrule
\multicolumn{4}{c}{\textbf{Ring $\mathbb{Z}_3$} ($|D_3|=6$, $B_3=5$, 3 Mosaics)} \\
\midrule
1 & 1 & $z_3$ & $1\omega^{1}$ \\
3 & 1 & $z_2 z_1$ & $1\omega^{3}$ \\
1 & 1 & $z_1^3$ & $1\omega^{1}$ \\
\midrule
\multicolumn{4}{c}{\textbf{Ring $\mathbb{Z}_4$} ($|D_4|=8$, $B_4=15$, 7 Mosaics)} \\
\midrule
1 & 1 & $z_4$ & $1\omega^{1}$ \\
4 & 1 & $z_3 z_1$ & $1\omega^{4}$ \\
3 & 2 & $z_2^2$ & $1\omega^{1} + 1\omega^{2}$ \\
6 & 2 & $z_2 z_1^2$ & $1\omega^{2} + 1\omega^{4}$ \\
1 & 1 & $z_1^4$ & $1\omega^{1}$ \\
\midrule
\multicolumn{4}{c}{\textbf{Ring $\mathbb{Z}_5$} ($|D_5|=10$, $B_5=52$, 12 Mosaics)} \\
\midrule
1 & 1 & $z_5$ & $1\omega^{1}$ \\
5 & 1 & $z_4 z_1$ & $1\omega^{5}$ \\
10 & 2 & $z_3 z_2$ & $2\omega^{5}$ \\
10 & 2 & $z_3 z_1^2$ & $2\omega^{5}$ \\
15 & 3 & $z_2^2 z_1$ & $3\omega^{5}$ \\
10 & 2 & $z_2 z_1^3$ & $2\omega^{5}$ \\
1 & 1 & $z_1^5$ & $1\omega^{1}$ \\
\midrule
\multicolumn{4}{c}{\textbf{Ring $\mathbb{Z}_6$} ($|D_6|=12$, $B_6=203$, 37 Mosaics)} \\
\midrule
1 & 1 & $z_6$ & $1\omega^{1}$ \\
6 & 1 & $z_5 z_1$ & $1\omega^{6}$ \\
15 & 3 & $z_4 z_2$ & $1\omega^{3} + 2\omega^{6}$ \\
10 & 3 & $z_3^2$ & $1\omega^{1} + 1\omega^{3} + 1\omega^{6}$ \\
15 & 3 & $z_4 z_1^2$ & $1\omega^{3} + 2\omega^{6}$ \\
60 & 6 & $z_3 z_2 z_1$ & $2\omega^{6} + 4\omega^{12}$ \\
15 & 5 & $z_2^3$ & $1\omega^{1} + 1\omega^{2} + 2\omega^{3} + 1\omega^{6}$ \\
20 & 3 & $z_3 z_1^3$ & $1\omega^{2} + 1\omega^{6} + 1\omega^{12}$ \\
45 & 8 & $z_2^2 z_1^2$ & $3\omega^{3} + 4\omega^{6} + 1\omega^{12}$ \\
15 & 3 & $z_2 z_1^4$ & $1\omega^{3} + 2\omega^{6}$ \\
1 & 1 & $z_1^6$ & $1\omega^{1}$ \\
\midrule
\multicolumn{4}{c}{\textbf{Ring $\mathbb{Z}_7$} ($|D_7|=14$, $B_7=877$, 93 Mosaics)} \\
\midrule
1 & 1 & $z_7$ & $1\omega^{1}$ \\
7 & 1 & $z_6 z_1$ & $1\omega^{7}$ \\
21 & 3 & $z_5 z_2$ & $3\omega^{7}$ \\
35 & 4 & $z_4 z_3$ & $3\omega^{7} + 1\omega^{14}$ \\
21 & 3 & $z_5 z_1^2$ & $3\omega^{7}$ \\
105 & 9 & $z_4 z_2 z_1$ & $3\omega^{7} + 6\omega^{14}$ \\
70 & 7 & $z_3^2 z_1$ & $4\omega^{7} + 3\omega^{14}$ \\
105 & 12 & $z_3 z_2^2$ & $9\omega^{7} + 3\omega^{14}$ \\
35 & 4 & $z_4 z_1^3$ & $3\omega^{7} + 1\omega^{14}$ \\
210 & 18 & $z_3 z_2 z_1^2$ & $6\omega^{7} + 12\omega^{14}$ \\
105 & 11 & $z_2^3 z_1$ & $7\omega^{7} + 4\omega^{14}$ \\
35 & 4 & $z_3 z_1^4$ & $3\omega^{7} + 1\omega^{14}$ \\
105 & 12 & $z_2^2 z_1^3$ & $9\omega^{7} + 3\omega^{14}$ \\
21 & 3 & $z_2 z_1^5$ & $3\omega^{7}$ \\
1 & 1 & $z_1^7$ & $1\omega^{1}$ \\
\midrule
\multicolumn{4}{c}{\textbf{Ring $\mathbb{Z}_8$} ($|D_8|=16$, $B_8=4140$, 354 Mosaics)} \\
\midrule
1 & 1 & $z_8$ & $1\omega^{1}$ \\
8 & 1 & $z_7 z_1$ & $1\omega^{8}$ \\
28 & 4 & $z_6 z_2$ & $1\omega^{4} + 3\omega^{8}$ \\
56 & 5 & $z_5 z_3$ & $3\omega^{8} + 2\omega^{16}$ \\
35 & 7 & $z_4^2$ & $1\omega^{1} + 1\omega^{2} + 2\omega^{4} + 3\omega^{8}$ \\
28 & 4 & $z_6 z_1^2$ & $1\omega^{4} + 3\omega^{8}$ \\
168 & 12 & $z_5 z_2 z_1$ & $3\omega^{8} + 9\omega^{16}$ \\
280 & 19 & $z_4 z_3 z_1$ & $3\omega^{8} + 16\omega^{16}$ \\
210 & 22 & $z_4 z_2^2$ & $1\omega^{2} + 4\omega^{4} + 10\omega^{8} + 7\omega^{16}$ \\
280 & 25 & $z_3^2 z_2$ & $2\omega^{4} + 12\omega^{8} + 11\omega^{16}$ \\
56 & 5 & $z_5 z_1^3$ & $3\omega^{8} + 2\omega^{16}$ \\
420 & 33 & $z_4 z_2 z_1^2$ & $1\omega^{4} + 12\omega^{8} + 20\omega^{16}$ \\
280 & 25 & $z_3^2 z_1^2$ & $2\omega^{4} + 12\omega^{8} + 11\omega^{16}$ \\
840 & 57 & $z_3 z_2^2 z_1$ & $9\omega^{8} + 48\omega^{16}$ \\
105 & 17 & $z_2^4$ & $1\omega^{1} + 2\omega^{2} + 5\omega^{4} + 8\omega^{8} + 1\omega^{16}$ \\
70 & 8 & $z_4 z_1^4$ & $1\omega^{2} + 1\omega^{4} + 4\omega^{8} + 2\omega^{16}$ \\
560 & 38 & $z_3 z_2 z_1^3$ & $6\omega^{8} + 32\omega^{16}$ \\
420 & 39 & $z_2^3 z_1^2$ & $3\omega^{4} + 21\omega^{8} + 15\omega^{16}$ \\
56 & 5 & $z_3 z_1^5$ & $3\omega^{8} + 2\omega^{16}$ \\
210 & 22 & $z_2^2 z_1^4$ & $1\omega^{2} + 4\omega^{4} + 10\omega^{8} + 7\omega^{16}$ \\
28 & 4 & $z_2 z_1^6$ & $1\omega^{4} + 3\omega^{8}$ \\
1 & 1 & $z_1^8$ & $1\omega^{1}$ \\
\bottomrule
\end{longtable}


\section{The Relational Datasets and Embedding Manifest}
\label{sec:relational-datasets}

To provide exhaustive machine-readable accessibility for computational mathematics, the complete analytical inventory across all $509$ mosaics has been structured into two normalised, greppable tab-separated tables housed in \path{data/}:

\begin{enumerate}
    \item \textbf{Master Mosaic Catalogue (\path{data/mosaics.tsv})}:
    Contains exactly $509$ records (one per dihedral orbit), cataloguing the global ID ($1 \dots 509$), ring order $n$, local ID ($1 \dots M_n$), shape $\lambda$, canonical slash representation, dihedral invariants ($|\mathcal{O}|$, $|\operatorname{Stab}|$, $s_{\text{sym}}$, chirality signature), chord dispersion energy $E$, polyatlas polygon profile, and the full relational \emph{block manifest}.
    
    \item \textbf{Relational Block Bridge (\path{data/mosaic_blocks.tsv})}:
    Contains exactly $1{,}983$ records (one per constituent block across all $509$ canonical mosaics), linking every individual block directly into the inscribed polygon census of \path{polyatlas/data/inngons.tsv}.
\end{enumerate}

\noindent\textbf{The Block Embedding Manifest}\\[0.4em]
While each mosaic $\mathcal{M} = \{B_1, \dots, B_s\}$ is dihedrally fixed as a collective set partition, its non-focus constituent blocks $B_j$ are generally translated by a spatial root offset and may be retrograded (reflected) relative to their isolated canonical presentation in \emph{The Polygonal Burnsidean}.

Every block embedding is encoded by a compact manifest token:
\begin{equation}
\label{eqn:block-manifest-syntax}
\text{[Class]}@\text{[Root]}[\text{Chirality}]:\text{[Realised Strides]}
\end{equation}
where:
\begin{itemize}
    \item \textbf{Class}: The canonical polygon class identifier $k\text{-}i$ defined in \texttt{polyatlas} (e.g.\ $3\text{-}1$ for the isosceles triangle on $\mathbb{Z}_4$, or $3\text{-}5$ for the chiral scalene triangle on $\mathbb{Z}_8$).
    \item \textbf{Root}: The starting vertex $v_0 = \min(B_j) \in \mathbb{Z}_n$ around the circular clock face.
    \item \textbf{Chirality}: Orientation sign relative to the canonical representative:
    \begin{itemize}
        \item \textbf{$+$ (Direct)}: A pure rotational translation of the canonical polygon.
        \item \textbf{$-$ (Retrograde)}: An enantiomorphic (reflected / retrograded) copy of the canonical polygon.
    \end{itemize}
    \item \textbf{Realised Strides}: The forward cyclic strides $(s_1, \dots, s_k)$ along the perimeter of block $B_j$ starting from $v_0$.
\end{itemize}

Concatenating these block tokens across the partition yields the self-contained mosaic manifest, as illustrated in Table~\ref{tab:sample-manifests}.

\begin{table}[htbp]
\centering
\footnotesize
\setlength{\tabcolsep}{2.8pt}
\caption{\label{tab:sample-manifests}Exemplar Mosaic Orbits across $\mathbb{Z}_1..\mathbb{Z}_8$ and their Relational Manifests}
\begin{tabular}{r c c l c l l}
\toprule
ID & Ring & Shape $\lambda$ & Canonical Mosaic & $|\mathcal{O}|$ & Polygon Profile & Relational Embedding Manifest \\
\midrule
1 & $\mathbb{Z}_1$ & $(1)$ & \texttt{/0/} & 1 & $1\text{-}1$ & \texttt{1-1@0+:1} \\
2 & $\mathbb{Z}_2$ & $(2)$ & \texttt{/01/} & 1 & $2\text{-}1$ & \texttt{2-1@0+:11} \\
3 & $\mathbb{Z}_2$ & $(1,1)$ & \texttt{/0/1/} & 1 & $1\text{-}1 + 1\text{-}1$ & \texttt{1-1@0+:2/1-1@1+:2} \\
4 & $\mathbb{Z}_3$ & $(3)$ & \texttt{/012/} & 1 & $3\text{-}1$ & \texttt{3-1@0+:111} \\
5 & $\mathbb{Z}_3$ & $(2,1)$ & \texttt{/0/12/} & 3 & $1\text{-}1 + 2\text{-}1$ & \texttt{1-1@0+:3/2-1@1+:12} \\
7 & $\mathbb{Z}_4$ & $(4)$ & \texttt{/0123/} & 1 & $4\text{-}1$ & \texttt{4-1@0+:1111} \\
9 & $\mathbb{Z}_4$ & $(2,2)$ & \texttt{/02/13/} & 1 & $2\text{-}2 + 2\text{-}2$ & \texttt{2-2@0+:22/2-2@1+:22} \\
10 & $\mathbb{Z}_4$ & $(2,2)$ & \texttt{/01/23/} & 2 & $2\text{-}1 + 2\text{-}1$ & \texttt{2-1@0+:13/2-1@2+:13} \\
26 & $\mathbb{Z}_6$ & $(6)$ & \texttt{/012345/} & 1 & $6\text{-}1$ & \texttt{6-1@0+:111111} \\
31 & $\mathbb{Z}_6$ & $(3,3)$ & \texttt{/024/135/} & 1 & $3\text{-}2 + 3\text{-}2$ & \texttt{3-2@0+:222/3-2@1+:222} \\
33 & $\mathbb{Z}_6$ & $(3,3)$ & \texttt{/013/245/} & 6 & $3\text{-}3 + 3\text{-}3$ & \texttt{3-3@0+:123/3-3@2-:213} \\
156 & $\mathbb{Z}_8$ & $(8)$ & \texttt{/01234567/} & 1 & $8\text{-}1$ & \texttt{8-1@0+:11111111} \\
168 & $\mathbb{Z}_8$ & $(4,4)$ & \texttt{/0246/1357/} & 1 & $4\text{-}6 + 4\text{-}6$ & \texttt{4-6@0+:2222/4-6@1+:2222} \\
382 & $\mathbb{Z}_8$ & $(2^4)$ & \texttt{/01/25/37/46/} & 16 & $2\text{-}1{+}2\text{-}2{+}2\text{-}3{+}2\text{-}4$ & \texttt{2-1@0+:17/2-3@2+:35/2-4@3+:44/2-2@4+:26} \\
\bottomrule
\end{tabular}
\end{table}

In Mosaic 33 on $\mathbb{Z}_6$, for example, the two chiral triangles partition the hexagon as an enantiomorphic direct-and-reflected pair (\texttt{3-3@0+:123} and \texttt{3-3@2-:213}). In Mosaic 382 on $\mathbb{Z}_8$, the four digon chords span four distinct circular distance classes ($1, 4, 3, 2$), realising the cyclic Skolem sequence through four distinct polyatlas classes ($2\text{-}1, 2\text{-}2, 2\text{-}3, 2\text{-}4$).

The complete visual plates displaying all $509$ mosaics follow in Chapter~\ref{chap:visual-atlas}.
